\documentclass[10pt,a4paper]{amsart}
\usepackage[margin=19mm]{geometry}
\usepackage{amsmath,amssymb,mathtools}
\usepackage[scr=rsfs,cal=euler]{mathalfa}
\usepackage{xfrac}
\usepackage[unicode,hidelinks,bookmarks=false]{hyperref}
\newcommand{\ie}{i.e.\ }
\newcommand{\cf}{cf.\ }
\newcommand{\resp}{respectively\ }
\newcommand{\Subsection}{\subsection}
\providecommand{\nobreakdash}{\nobreak\hbox{-}}
\setlength{\emergencystretch}{2em}
\multlinegap0pt
\usepackage[shortlabels]{enumitem}
\usepackage{amssymb}
\usepackage{bbm}
\usepackage{tikz}
\usepackage{booktabs}
\usepackage{multirow}
\usepackage{float}
\newtheorem{definition}{Definition}
\newtheorem{proposition}{Proposition}
\newcommand{\bF}{\mathbf{F}}
\newcommand{\bT}{\mathbf{T}}
\newcommand{\bU}{\mathbf{U}}
\newcommand{\ipl}{\ensuremath{\mathbf{IPL}}}
\title{Intuitionism with truth-tables: A Decision Procedure for \textbf{IPL} Based on RNmatrices}
\author{Renato Leme, Marcelo Coniglio, Bruno Lopes}
\date{Source: arXiv:2308.13664v4 (CC BY 4.0)}
\begin{document}
\maketitle

\noindent\emph{Source and licence.} Intuitionism with truth-tables: A Decision Procedure for \textbf{IPL} Based on RNmatrices. Version arXiv:2308.13664v4 (CC BY 4.0), distributed by arXiv under the Creative Commons Attribution 4.0 International licence, http://creativecommons.org/licenses/by/4.0/, as recorded in the arXiv licence metadata for this exact version and shown on the abstract page https://arxiv.org/abs/2308.13664v4. Full licence notice: \texttt{licenses/arxiv-2308.13664-CC-BY-4.0.txt}.\par\smallskip

\noindent\textit{Editorial scope.} Counted unit: the article's proposition seq-rules-sound (source0.tex lines 1007--1009). The article's complete original proof of it is delivered with it (source0.tex lines 1010--1041, one \texttt{proof} environment of 32 lines). No mathematics is written, repaired or re-derived: the statement and the proof are the article's own lines, in the article's own notation, and no retained line is substituted or re-flowed.  Retained uncounted as the source-local context the proof needs: lines 927--940.  Nothing was omitted from any retained span, so the package carries no omission record.  No retained line contains a citation, so the wrapper carries no bibliography and no bibliography entry.  No retained line contains a figure environment, an image-inclusion command or drawing code, so no image is delivered and the package declares no asset dependency.  Editorial formatting only, in addition: an AMS \texttt{amsart} preamble that restates the article's own declarations and macros used by the retained lines, namely the environments \texttt{definition}, \texttt{proposition} on separate counters, together with the standard packages those lines need.  The retained environments print in this document as Definition 1, Proposition 1 in order of appearance, whereas the article prints the same items with the numbering of its own full document.  The complete native source of the article as distributed in the arXiv e-print for this exact version is bundled unchanged as \texttt{sources/146118/source0.tex}.
\begin{definition}\label{def:level_valuations_IPL} (Level valuations for \ipl)
The set of $n$th-level valuations $\mathcal{L}^{IPL}_n$ for \ipl, where $n \in \omega$, is defined as follows: 

    \begin{enumerate}[label=\roman*)]
        \item $\mathcal{L}_0^{IPL} = Val^{nd}_+(\mathcal{M}_{IPL})$;
        \item $\mathcal{L}^{IPL}_{n + 1} = \{ v \in \mathcal{L}^{IPL}_n \mid \forall \alpha \in For(\Omega)$, if $v(\alpha)=\bU$ then there exists $w \in \mathcal{L}^{IPL}_n$ such that $w(\alpha)=\bF$ and,  for every $\beta \in  For(\Omega)$: $v(\beta)=\bT$ implies $w(\beta)=\bT\}$.
    \end{enumerate}
    The set of level valuations in $\mathcal{M}_{IPL}$ is defined as the intersection of the sets of $n$th-level valuations, for $n\geq0$:

    \[
        \mathcal{L}_{IPL} = \bigcap^\infty_{n \geq 0} \mathcal{L}^{IPL}_n
    \]

\end{definition}

\begin{proposition} \label{seq-rules-sound}
For every $k \geq 0$ and every sequent $\Gamma \Rightarrow\Delta$: if $\Gamma \vdash_{IPL}^k \Delta$, then $\Gamma\models_k\Delta$.
\end{proposition}

\begin{proof}
By induction on $k$.\\[1mm]
%
{\bf Base} $k=0$: Suppose that $\Gamma \vdash_{IPL}^0 \Delta$. Then  $\Gamma \Rightarrow\Delta$ is an instance of the axiom ($Ax$), that is:  $\Gamma \Rightarrow\Delta = \alpha,\Gamma' \Rightarrow \alpha$ for some multiset $\Gamma'$. Let $v \in \mathcal{L}_0^{IPL}$ such that $v \models \Gamma' \cup \{\alpha\}$. In particular, $v \models \alpha$, hence $v \models \Gamma \Rightarrow\Delta$. From this,  $\Gamma \models_0 \Delta$.\\[1mm]
%
{\bf Inductive step}: Suppose that, for any sequent  $\Gamma \Rightarrow\Delta$, if $\Gamma \vdash_{IPL}^j \Delta$ then  $\Gamma\models_j\Delta$, for every $j \leq k$, for a given $k \geq 0$ (Induction Hypothesis, IH). Let $\Gamma \Rightarrow\Delta$ be a sequent such that $\Gamma \vdash_{IPL}^{k+1} \Delta$. If $\Gamma \Rightarrow \Delta$ is an instance of the axiom, it can be proven as above that  $\Gamma\models_{k+1}\Delta$. Suppose now that $\Gamma \Rightarrow\Delta$ was obtained from other sequents by applying a rule of {\bf LI} in a derivation of depth $k+1$. We have the following cases:\\[1mm]
%
($contr\Rightarrow$): Suppose that  $\Gamma \Rightarrow\Delta= \alpha,\Gamma' \Rightarrow\Delta$ is obtained from $\alpha,\alpha,\Gamma'\Rightarrow\Delta$ by applying rule ($contr\Rightarrow$). Then $\alpha,\alpha,\Gamma' \vdash_{IPL}^k \Delta$ and so $\alpha,\alpha,\Gamma' \models_k \Delta$, by (IH). But the latter clearly implies that $\alpha,\Gamma' \models_k \Delta$ and so  $\alpha,\Gamma' \models_{k+1} \Delta$,  given that $\mathcal{L}_{k+1}^{IPL} \subseteq \mathcal{L}_k^{IPL}$.\\[1mm]
%
($\Rightarrow w$): Suppose that  $\Gamma \Rightarrow\Delta= \Gamma \Rightarrow\alpha$ is obtained from $\Gamma\Rightarrow$ by ($\Rightarrow w$). Then  $\Gamma \vdash_{IPL}^k $ and so, by (IH), $\Gamma\models_k$.  This means that, for every $v \in \mathcal{L}_k^{IPL}$, $v\not\models \Gamma$ and so $v\models\Gamma\Rightarrow \alpha$. From this it follows that $\Gamma\models_{k+1}\Delta$,  given that $\mathcal{L}_{k+1}^{IPL} \subseteq \mathcal{L}_k^{IPL}$.\\[1mm]
%
($cut$): Assume that $\Gamma \Rightarrow\Delta$ follows from $\Gamma\Rightarrow\alpha$ and $\alpha, \Gamma\Rightarrow\Delta$ by ($cut$) rule. Then  $\Gamma\vdash_{IPL}^k\alpha$ and $\alpha, \Gamma\vdash_{IPL}^k\Delta$ and so, by (IH),  $\Gamma \models_k\alpha$ and $\alpha, \Gamma \models_k\Delta$.
Suppose first that $\Delta=\emptyset$.  Let $v \in \mathcal{L}_{k+1}^{IPL}$ such that $v \models\Gamma$; in particular, $v \in \mathcal{L}_{k}^{IPL}$ and so $v \models\Gamma\Rightarrow\alpha$. From this it follows that  $v\models\alpha$. But this is impossible, given that $v \models \Gamma \cup\{\alpha\} \Rightarrow$, by (IH). From this, $v \not\models\Gamma$ for every $v\in \mathcal{L}_{k+1}^{IPL}$. That is, $v\models\Gamma\Rightarrow$, for every $v\in \mathcal{L}_{k+1}^{IPL}$. Now, suppose that $\Delta=\{\gamma\}$ for some formula $\gamma$, and let $v\in \mathcal{L}_{k+1}^{IPL}$ such that $v \models \Gamma$. Once again, since $v \models\Gamma\Rightarrow\alpha$ we infer that  $v\models\alpha$. But then $v \models \Gamma \cup \{\alpha\}$ and so $v\models\gamma$, given that $v \models\Gamma \cup\{\alpha\}\Rightarrow \gamma$. This means that $v \models \Gamma\Rightarrow\Delta$, for every $v\in \mathcal{L}_{k+1}^{IPL}$. That is, $\Gamma\models_{k+1}\Delta$.\\[1mm]
%
($\vee\Rightarrow$): Assume that  $\Gamma \Rightarrow\Delta=\alpha \vee\beta,\Gamma' \Rightarrow\Delta$ follows from $\alpha, \Gamma'\Rightarrow\Delta$ and $\beta, \Gamma'\Rightarrow\Delta$ by ($\vee\Rightarrow$). By (IH),  $\alpha, \Gamma' \models_k \Delta$ and $\beta, \Gamma' \models_k \Delta$. Suppose first that $\Delta=\emptyset$. By hypothesis, if $v \in \mathcal{L}_{k+1}^{IPL}$ is such that $v\models\Gamma'$, then $v\not\models\alpha$ and $v\not\models\beta$, and so $v\not\models\alpha \vee\beta$. This means that $v \models \Gamma' \cup \{\alpha\vee\beta\}\Rightarrow$, for every $v\in \mathcal{L}_{k+1}^{IPL}$. Now, if $\Delta=\{\gamma\}$ for some formula $\gamma$, let $v\in \mathcal{L}_{k+1}^{IPL}$. By hypothesis, if  $v\models\Gamma'$ and $v\models\alpha$, then $v\models\gamma$; and if  $v\models\Gamma'$ and $v\models\beta$, then $v\models\gamma$. From this, if  $v\models\Gamma'$ and $v\models\alpha\vee\beta$ then, by definition of $\mathcal{M}_{IPL}$, either $v\models\alpha$ or $v\models\beta$  and so, by hypothesis,  $v\models\gamma$. That is, $v \models \Gamma\Rightarrow\Delta$ for every $v\in \mathcal{L}_{k+1}^{IPL}$ and so  $\Gamma\models_{k+1}\Delta$.\\[1mm]
%
($\Rightarrow\vee_i$): It is enough to observe that, if either $v\models \alpha$ or $v\models \beta$ then $v\models \alpha\vee\beta$, by definition of $\mathcal{M}_{IPL}$.\\[1mm]
%
($\wedge\Rightarrow$): Suppose that $\Gamma \Rightarrow\Delta=\alpha \land \beta,\Gamma' \Rightarrow\Delta$ follows from $\alpha, \beta,\Gamma'\Rightarrow\Delta$ by ($\wedge\Rightarrow$). Assume first that $\Delta=\emptyset$. Then, if $v \in \mathcal{L}_{k+1}^{IPL}$ such that $v \models\Gamma'$  it follows that $v \not\models\{\alpha,\beta\}$, since $\mathcal{L}_{k+1}^{IPL} \subseteq \mathcal{L}_k^{IPL}$ and $\alpha, \beta,\Gamma'\models_{k}\Delta$ by (IH). From this, either  $v \not\models\alpha$ or  $v \not\models\beta$. By definition of $\mathcal{M}_{IPL}$,  $v \not\models\alpha \land\beta$. That is, $v \not\models\Gamma' \cup \{\alpha\land\beta\}$ for every $v\in \mathcal{L}_{k+1}^{IPL}$, i.e., $\alpha\wedge\beta, \Gamma'\models_{k+1}$. In turn, if  $\Delta=\{\gamma\}$ for some formula $\gamma$, let $v \in \mathcal{L}_{k+1}^{IPL}$ such that $v \models\Gamma' \cup \{\alpha\land\beta\}$. By definition of $\mathcal{M}_{IPL}$, $v \models\Gamma' \cup \{\alpha,\beta\}$ and so $v\models\gamma$, by (IH) and the fact that  $v \in \mathcal{L}_k^{IPL}$. That is, $v\models \Gamma' \cup \{\alpha \land\beta\}\Rightarrow\Delta$ for every $v\in \mathcal{L}_{k+1}^{IPL}$.\\[1mm]
%
($\Rightarrow\wedge$): It is enough to observe that, if both $v\models \alpha$ and $v\models \beta$ then $v\models \alpha\wedge\beta$, by definition of $\mathcal{M}_{IPL}$.\\[1mm]
%
($\to\Rightarrow$) Suppose that $\Gamma \Rightarrow\Delta= \alpha \to \beta,\Gamma'\Rightarrow \Delta$ is obtained from $\Gamma' \Rightarrow \alpha$ and $\beta, \Gamma' \Rightarrow\Delta$ by ($\to\Rightarrow$). By (IH),  $\Gamma'\models_{k} \alpha$ and $\beta, \Gamma' \models_{k}\Delta$. Assume first that $\Delta=\emptyset$. Let $v \in \mathcal{L}_{k+1}^{IPL}$ such that $v \models\Gamma' \cup\{\alpha \to \beta\}$. Then, $v  \in \mathcal{L}_{k}^{IPL}$ such that $v \models\Gamma'$ and so $v\models \alpha$, by hypothesis. Since $v \models\alpha \to \beta$ then  $v \models \beta$, by definition of  $\mathcal{M}_{IPL}$. That is, $v  \in \mathcal{L}_{k}^{IPL}$ is such that $v \models\Gamma'\cup \{\beta\}$, which contradicts the fact that $\beta, \Gamma' \models_{k}$. From this,  $v \not\models\Gamma' \cup\{\alpha \to \beta\}$ for every $v \in \mathcal{L}_{k+1}^{IPL}$, which means that $\alpha \to\beta,\Gamma'\models_{k+1}$. Now, suppose that $\Delta=\{\gamma\}$ for some formula $\gamma$.  Let $v \in \mathcal{L}_{k+1}^{IPL}$ such that $v \models\Gamma' \cup\{\alpha \to \beta\}$. By reasoning as above,   $v  \in \mathcal{L}_{k}^{IPL}$ is such that $v \models\Gamma'\cup \{\beta\}$. But then $v \models \gamma$, by hypothesis. This implies that $v \models\Gamma' \cup\{\alpha \to \beta\} \Rightarrow \gamma$ for every $v \in \mathcal{L}_{k+1}^{IPL}$. \\[1mm]
%
($\Rightarrow\to$): Assume that $\Gamma \Rightarrow\Delta= \Gamma\Rightarrow \alpha \to \beta$ is obtained from $\alpha,  \Gamma \Rightarrow \beta$ by ($\Rightarrow\to$). By (IH),  $\alpha,  \Gamma\models_{k} \beta$. Let $v \in \mathcal{L}_{k+1}^{IPL}$ such that $v \models \Gamma$. Suppose that $v(\alpha \to \beta)=\bF$. By definition of $\mathcal{M}_{IPL}$, $v(\alpha)=\bT$ and $v(\beta) \in \{\bU,\bF\}$. That is, $v \models\Gamma \cup\{\alpha\}$ but $v \not\models\beta$, which contradicts the fact that $\alpha,  \Gamma\models_{k} \beta$, given that $v \in \mathcal{L}_{k}^{IPL}$. Hence, $v(\alpha \to \beta)\neq\bF$. Suppose now that $v(\alpha \to \beta)=\bU$. By Definition~\ref{def:level_valuations_IPL} of level valuations, there exists $w \in \mathcal{L}_{k}^{IPL}$ such that  $w(\alpha \to \beta)=\bF$ and $w \models \Gamma$. By reasoning as above, $w(\alpha)=\bT$ and $w(\beta) \in \{\bU,\bF\}$. But then $w \models \Gamma \cup\{\alpha\}$ and $w \not\models\beta$, which contradicts the fact that  $\alpha,  \Gamma\models_{k} \beta$. From this, it follows that  $v(\alpha \to \beta)=\bT$. That is, $v\models \Gamma\Rightarrow \alpha\to\beta$ for every $v \in \mathcal{L}_{k+1}^{IPL}$.\\[1mm]
%
($\neg\Rightarrow$): Assume that $\Gamma \Rightarrow\Delta= \neg \alpha,\Gamma'\Rightarrow$ is obtained from $\Gamma' \Rightarrow \alpha$ by ($\neg\Rightarrow$). By (IH),  $\Gamma'\models_{k} \alpha$. Let $v \in \mathcal{L}_{k+1}^{IPL}$ such that $v \models \Gamma' \cup\{\neg\alpha\}$. Then, $v \in \mathcal{L}_{k}^{IPL}$ such that $v \models \Gamma'$ and so $v \models \alpha$, by hypothesis, which contradicts the fact that $v \models \neg\alpha$.  From this,  $v \not\models\Gamma' \cup\{\neg\alpha\}$ for every $v \in \mathcal{L}_{k+1}^{IPL}$, which means that $\neg\alpha,\Gamma'\models_{k+1}$. \\[1mm]
%
($\Rightarrow\neg$): Assume that $\Gamma \Rightarrow\Delta= \Gamma\Rightarrow \neg\alpha$ is obtained from $\alpha,  \Gamma \Rightarrow$ by ($\Rightarrow\neg$). By (IH),  $\alpha,  \Gamma\models_{k} $. Let $v \in \mathcal{L}_{k+1}^{IPL}$ such that $v \models \Gamma$. Suppose that $v(\neg\alpha)=\bF$. By definition of $\mathcal{M}_{IPL}$, $v(\alpha)=\bT$, hence $v \models\Gamma \cup\{\alpha\}$, which contradicts the fact that $\alpha,  \Gamma\models_{k}$, since $v \in \mathcal{L}_{k}$. Now, suppose that $v(\neg\alpha)=\bU$. By Definition~\ref{def:level_valuations_IPL} of level valuations, there exists $w \in \mathcal{L}_{k}^{IPL}$ such that  $w(\neg\alpha)=\bF$ and $w \models \Gamma$. As proven above, $w(\alpha)=\bT$ and so $w \models \Gamma \cup\{\alpha\}$, contradicting the fact that  $\alpha,  \Gamma\models_{k}$. From this, it follows that  $v(\neg\alpha)=\bT$. That is, $v\models \Gamma\Rightarrow \neg\alpha$ for every $v \in \mathcal{L}_{k+1}^{IPL}$.

This concludes the proof.
\end{proof}

\end{document}
